\section{Proper escape and arithmetic return laws}\label{sec:levels}

\subsection{Bounded half-orbits and two-sided escape}

The passage from periodic classification to integral escape uses
discreteness rather than a real dynamical estimate.

\begin{lemma}[A discrete-bijection principle]\label{lem:discrete}
Let $F$ be a bijection of a discrete subset $X$ of a finite-dimensional
real vector space such that every bounded subset of $X$ is finite.
A point of $X$ has a bounded forward or backward orbit if and only
if it is periodic. Every nonperiodic two-sided orbit visits each bounded
subset only finitely many times.
\end{lemma}
\begin{proof}
A bounded forward orbit lies in a finite set, so $F^iP=F^jP$
for some $0\leq i<j$. Applying $F^{-i}$ gives
$P=F^{j-i}P$. The same argument with $F^{-1}$ treats backward
orbits. Conversely, a periodic orbit is finite and bounded.

For a nonperiodic point, the points $F^nP$, $n\in\Z$, are
pairwise distinct: any equality would give a period by applying an
inverse iterate. A finite subset can contain only finitely many of
these distinct points. This proves the last assertion.
\end{proof}

Apply the lemma to $X=\Z^3$ and $F=T$. For any finite radius
$R$, only finitely many integers $n$ have
$T^nP\in[-R,R]^3$ when $P$ is nonperiodic. Thus, beyond the
last visit in either time direction, $\|T^nP\|_\infty>R$.
Since $R$ is arbitrary, this proves \eqref{eq:escape} and completes
the boundedness part of Theorem~\ref{thm:main}.

\subsection{Every level and every ordinary time}

For $k\in\Z$, define
\begin{equation}\label{eq:indicators}
 s_k=\ind_{\{k=m^2\text{ for some }m\geq1\}},\qquad
 q_k=\ind_{\{k=m^2-m+2\text{ for some }m\geq1\}}.
\end{equation}
Both parametrizations are injective on positive integers. The second
condition is equivalently that $4k-7$ be a positive odd square,
because $4(m^2-m+2)-7=(2m-1)^2$.
Let $c_d(k)$ denote the number of orbits of least period $d$ on
$S_k(\Z)$. The classification immediately gives
\begin{align}
 c_1(k)&=\ind_{k=0}+\ind_{k=4},&
 c_3(k)&=\ind_{k=4},\nonumber\\
 c_4(k)&=q_k,& c_6(k)&=s_k,& c_{12}(k)&=q_k,
 \label{eq:orbit-counts}
\end{align}
and $c_d(k)=0$ for every other $d$. In particular these counts
are finite on each fixed level, whether or not that level contains
infinitely many nonperiodic integral points.

\begin{corollary}[All-time fixed-point counts]\label{cor:fixed}
For every $k\in\Z$ and $n\geq1$,
\begin{align}
 F_k(n):=\#\Fix(T^n;S_k(\Z))
 &=\ind_{k=0}+\ind_{k=4}\bigl(1+3\ind_{3\mid n}\bigr)
 \nonumber\\
 &\quad+6s_k\ind_{6\mid n}
 +q_k\bigl(4\ind_{4\mid n}+12\ind_{12\mid n}\bigr).
 \label{eq:fixed-counts}
\end{align}
\end{corollary}
\begin{proof}
An orbit of least period $d$ contributes all $d$ of its points
to $\Fix(T^n)$ if $d\mid n$, and none otherwise. Therefore
$F_k(n)=\sum_{d\mid n}dc_d(k)$, and
\eqref{eq:orbit-counts} gives the result.
\end{proof}

\begin{definition}[Fixed-level dynamical zeta function]
The ordinary dynamical zeta function of the integral points on a
fixed level is the formal power series
\begin{equation}\label{eq:zeta-definition}
 \zeta_k(t)=\exp\left(\sum_{n\geq1}\frac{F_k(n)}{n}t^n\right)
 \in 1+t\Q[[t]].
\end{equation}
\end{definition}

\begin{corollary}[Exact fixed-level zeta function]\label{cor:zeta}
For every integer $k$,
\begin{equation}\label{eq:zeta}
 \zeta_k(t)=
 \frac{1}{
 (1-t)^{\ind_{k=0}+\ind_{k=4}}
 (1-t^3)^{\ind_{k=4}}
 (1-t^6)^{s_k}
 (1-t^4)^{q_k}
 (1-t^{12})^{q_k}}.
\end{equation}
This is also an analytic identity for $|t|<1$ and gives a rational
continuation to the complex plane.
\end{corollary}
\begin{proof}
For a single orbit of length $d$, its contribution to the exponent
of \eqref{eq:zeta-definition} is
\[
 \sum_{r\geq1}\frac{d}{dr}t^{dr}
 =-\log(1-t^d).
\]
Only finitely many orbits occur at any fixed $k$, so summing these
contributions gives
$\zeta_k(t)=\prod_d(1-t^d)^{-c_d(k)}$, which is
\eqref{eq:zeta}. Formula \eqref{eq:fixed-counts} bounds $F_k(n)$
uniformly in $n$, so the defining logarithmic series is absolutely
convergent for $|t|<1$ as well.
\end{proof}

This product is over source periodic orbits. Its index is a period,
not a rational prime or a residue-field degree. The rationality in
\eqref{eq:zeta} follows from the finite periodic locus on the
fixed level, not from an additional arithmetic correspondence.

\subsection{The intersection of the two supports}

\begin{lemma}\label{lem:intersection}
The sets $\{r^2:r\geq1\}$ and
$\{m^2-m+2:m\geq1\}$ intersect only at $4$.
\end{lemma}
\begin{proof}
Suppose $r^2=m^2-m+2$ with $r,m\geq1$. Completing the square
and factoring give
\begin{equation}\label{eq:factor-seven}
 (2r-2m+1)(2r+2m-1)=7.
\end{equation}
Since $r^2=(m-\tfrac12)^2+\tfrac74$, we have
$r>m-\tfrac12$. Both factors in \eqref{eq:factor-seven} are
therefore positive integers. The second is larger than the first
because their difference is $4m-2>0$. As $7$ is prime, they
must be $1$ and $7$. Adding the two equations gives $r=2$,
and subtracting gives $m=2$. Conversely those values give level $4$.
\end{proof}

The total periodic-point count is
\begin{equation}\label{eq:total}
 \sum_{d\geq1}dc_d(k)=
 \ind_{k=0}+4\ind_{k=4}+6s_k+16q_k.
\end{equation}
Lemma~\ref{lem:intersection} makes its alternatives disjoint, as
recorded in Table~\ref{tab:levels}; this proves
Corollary~\ref{cor:support}.

\begin{table}[htbp]
\centering
\begin{tabular}{@{}lll@{}}
\toprule
Level condition & Periodic orbits & Number of points\\
\midrule
$k=0$ & $O$ & $1$\\
$k=4$ & $E,D,A_2,B_2,C_2$ & $26$\\
$k=m^2\ne4$, $m\geq1$ & $A_m$ & $6$\\
$k=m^2-m+2\ne4$, $m\geq1$ & $B_m,C_m$ & $16$\\
All remaining $k\in\Z$ & None & $0$\\
\bottomrule
\end{tabular}
\caption{The complete periodic-point count on every integer level.
The alternatives are mutually exclusive by Lemma~\ref{lem:intersection}.}
\label{tab:levels}
\end{table}

There is no ordinary finite-count zeta function on the unrestricted
lattice obtained by removing $k$ from the notation. For example,
every point of $A_m$ is fixed by $T^6$, and the disjoint axis
orbits exist for all $m\geq1$. Hence
$\#\Fix(T^6;\Z^3)=\infty$. The fixed-level quantifier in
\eqref{eq:zeta-definition} is indispensable.
